% =============================================================
%  HeInSulfides_abstract.tex
%  Shared abstract text for Chapter 3 (sulfide He step heats and
%  geochronology). Kept separate from any \section{Abstract}/quote
%  wrapper for the same reason as Chapter 1 and 2's abstract files:
%  the surrounding heading/environment is not class-agnostic, but
%  this raw text is.
%
%  \input{} by Chapter3/chapter3.tex (dissertation chapter).
%  Edit this file, not chapter3.tex, to change the abstract text.
% =============================================================
Sulfide mineralization is a near-ubiquitous process in the formation of
ore deposits of valuable chalcophile and siderophile elements. Helium
isotope studies of sulfides have been applied to infer the origins of
mineralizing hydrothermal fluids, and some pioneering studies have
applied the U-Th/He method to date sulfides, especially pyrite
\citep{yakubovich2020,yakubovich2021}. Yet, the interpretation of helium
isotope data or U-Th/He ages requires an understanding of the conditions
under which helium in sulfides behaves as an open or closed system, as
well as the mechanisms of helium loss, a process which is poorly
understood in sulfides. Furthermore, sulfides systematically reject U-Th We explore the systematics of U-Th/He in
sulfides through two lenses: 1) an example application of the sulfide
U-Th/He geochronometer to two relatively successful test cases, and 2)
helium step heat experiments of sulfide phases. We report
pyrite/arsenopyrite U-Th/He ages for Au-ores in the Gympie formation of
Queensland Australia, which yield a mean age of 253 $\pm$ 21 Ma, and span the gap between
reported Ar/Ar and K/Ar ages of the the country rock and the mineralized quartz veins \citep{holcombe1997, little1995}.
and the 2.7 Ga Kidd Creek VMS deposit of Ontario, Canada, which exhibits variable helium retention between an
unconstrained age older than 1.5 Ga and a cooling age \textasciitilde{}
500 Ma, similar to but older than the regional apatite cooling ages reported by
\citet{sturrock2024}.

Helium step heat experiments in sulfides are challenging because
sulfides readily react/decompose in the high vacuum, high temperature
conditions of the experiments, producing results that reflect phase
decomposition/reaction kinetics that are not geologically relevant. We
address this challenge by performing helium step heat experiments at
temperatures up to 600\degC{} on natural single-crystal pyrite, arsenopyrite,
and galena implanted with uniform 3He through proton bombardment sealed
in helium-permeable silica glass tubes stabilized under high sulfur (±
arsenic) fugacity to prevent phase decomposition during analysis. We
were unsuccessful in extending this method to stabilize chalcopyrite,
chalcocite, or marcasite during step heat experiments: the helium
retentivity of these phases in nature remains unclear.
Retrograde-prograde thermal cycles on sulfur-stabilized galena support a
Multiple Diffusion Domain (MDD) helium loss mechanism with an activation
energy around \ChThreeCCpEa\ kJ/mol -- the chalcopyrite step-heat
regression (\ChThreeCCpN\ measurements, R\textsuperscript{2}=\ChThreeCCpRtwo)
used as a shared reference line for both phases -- with retentivities
generally similar to apatite. Comparisons of helium step heat data of
sulfur-stabilized pristine and deformed pyrite point towards a multipath
diffusion model, in which dislocations and fine fractures serve as high
activation energy (\ChThreeRustEa\ $\pm$ \ChThreeRustEaErr\ kJ/mol) fast
pathways in deformed pyrites with bulk volume diffusion following a
lower activation energy of \ChThreeLVCOEa\ $\pm$ \ChThreeLVCOEaErr\
kJ/mol. Helium step heat data of arsenopyrite are complex and may be
complicated by exchange of sulfur and arsenic with the vapor phase, but
qualitatively demonstrate very high helium retentivity. These
experiments demonstrate that pyrite and arsenopyrite are similarly
helium-retentive as monzonite-structure phosphates and are more
retentive of helium than natural zircon, apatite, hematite, or goethite.
